\chapter{Upgrading Boogie}
\label{sec:upgrading-boogie}

The first contribution of this work was updating the Boogie version underlying the proof generation repository.
The last official Boogie commit merged into this repository was dated February 4th, 2022, corresponding to Boogie version 2.11.4.\footnote{\url{https://github.com/viperproject/boogie-proofgen/commit/0b9f14d15b757c2b5e3ca5d9fa84a5830e9c00bf}}
At the onset of this project, the newest available release was Boogie version 3.5.5 (released August 7th, 2025).
Given the elapsed time and the major version jump, substantial changes to the Boogie codebase were expected.

\section{Default Type Encoding Change}
\label{sec:default-type-encoding}

One highly relevant change between these versions is that the default type encoding in Boogie shifted from predicate to monomorphic.
The proof generation module currently provides better support for predicate type encoding,
as explained later in Section~\ref{sec:type-encoding}.
Consequently, it is strictly necessary to explicitly pass the predicate option flag (\lstinline[columns=fixed]{/typeEncoding:p})
to ensure the tool exhibits the same behavior as the older version.

\section{Checking SMT Changes}

Before delving directly into the code update, we first investigated whether the underlying SMT encoding had significantly changed.
Because proof generation relies heavily on the structure of the SMT output, structural changes to the verification conditions (VC) could break the module.
We compared the SMT output generated by Boogie v2.11.4 against that of Boogie v3.5.5.
% [relevant to compare proofgen?]
% mention? disable control flow VC

We generated the SMT output for all tests in the \texttt{ProofGenerationBenchmarks} folder (roughly 130 tests).
% TODO ref
While every file exhibited differences, a manual inspection of a subset revealed that the vast majority of changes were structurally insignificant.
Most files only showed variations in the inclusion or ordering of SMT options, rather than changes to the actual VC logic.
Because these options assist the SMT solver's internal satisfiability heuristics, they fall outside the concern of the proof generation logic.

However, a few files did exhibit changes to the VC.
These mostly concerned predicate type axioms, such as:

% \[
% \forall x:U. \operatorname{type}\ x = \operatorname{intType}
% \implies \operatorname{int\_2\_U} (\operatorname{U\_2\_int}\ x) = x
% \]


\begin{lstlisting}[language=SMT, numbers=none,
caption={SMT type axiom excerpt}]
forall ((x T@U) )
  (=> (= (type x) intType) (= (int_2_U (U_2_int x)) x))
\end{lstlisting}

This specific axiom asserts that converting an integer to a boxed value and back preserves its identity.
In older Boogie versions, these axioms were generated for \emph{every} type, whereas newer versions only generate them for types explicitly used in the input program.
Furthermore, Boogie dropped type ordering entirely,\footnote{\url{https://github.com/boogie-org/boogie/pull/667/}}
which is advantageous, as these orderings were previously deleted manually by the proof generation module.
Both of these changes appeared compatible with the existing Proofgen architecture.
\footnote{Later (in Section~\ref{sec:smt-change}), we discovered that four test cases did contain a minor but significant structural change affecting proof generation. Because all files had textual differences and the affected cases also contained polymorphic type axioms, this specific logical difference was difficult to spot in advance.}

\section{Merging Process}

Having observed no immediately critical changes in the SMT encoding, we proceeded to integrate Proofgen with the newest version of Boogie.
Because Boogie had undergone extensive refactoring, and because Proofgen was originally designed to be as minimally intrusive to the existing codebase as possible~\cite{parthasarathy2021formally},
we decided to manually apply the diff between Proofgen and Boogie v2.11.4 onto Boogie v3.5.5.

The generated diff contained fewer than 700 lines of changes spread across 15 files,
whereas the diff between Boogie v2.11.4 and v3.5.5 impacted over 300 files.
Therefore, we forked the official Boogie repository, created a new branch based on v3.5.5,
and manually ported the code file by file until every Proofgen modification was successfully integrated and compiling.
% mention periodic checking of smt changes?

\section{Preserving Git History}

A potential downside to this manual porting approach is the loss of the original Proofgen git history,
as the ported code would be attributed entirely to new commits.
We mitigated this through the following Git workflow:
First, we forked the original Proofgen repository and created a new working branch.
We then added our newly ported v3.5.5 repository as a remote, fetched it, and merged it into the working branch.
This merge incorporated all the architectural updates from v3.5.5 while preserving the legacy Proofgen history.
After resolving the merge conflicts, we verified that the history from both Boogie and Proofgen remained largely intact;
only our explicit adaptations, such as renaming and moving specific code blocks, were attributed to the new commits.

\section{Implementation Changes}

The following sections detail the specific code adjustments required to successfully complete the merge.

\subsection{Refactoring Changes}

Many of the required changes were standard software upkeep.
Several variables required renaming (e.g., changing \texttt{if.thn} to \texttt{if.Thn})
as fields were made private and exposed via public C\# properties.\footnote{\url{https://learn.microsoft.com/en-us/dotnet/csharp/programming-guide/classes-and-structs/properties}}
Files, methods, and individual lines of code had been moved, split, or merged across the Boogie repository,
necessitating corresponding relocations of the Proofgen modifications.
Additionally, some interfaces had evolved, requiring adjustments to the Proofgen code, while certain utility functions were rendered obsolete.

\subsection{CLI Changes}

A more involved architectural change was the removal of Boogie's static command-line options variable.
Previously, user-defined command-line arguments could be globally accessed via the static variable \lstinline[columns=fixed]{CommandLineOptions.Clo}.
% TODO line break
\footnote{\url{https://github.com/boogie-org/boogie/pull/501}}
We adopted the developers' recommended approach and introduced our own static variable, similar to how Dafny initially handled the transition.\footnote{\url{https://github.com/dafny-lang/dafny/pull/1866}}
While Dafny subsequently refactored this into a non-static implementation,\footnote{\url{https://github.com/dafny-lang/dafny/pull/3663}}
doing so in Proofgen would require pervasive code changes,
so we have deferred this architectural modernization to future work.

\subsection{SMT Changes and Proof Fixes}
\label{sec:smt-change}

After resolving all compilation errors, we ran the proof generation module against the benchmark suite to ensure parity with the older version.
The test run yielded 9 failures. Five of these were known failures inherited from before the merge, leaving 4 new regressions.

By analyzing the SMT output of one of the failing tests, we identified a structural difference we had previously missed:

% word diff colouring?
\begin{lstlisting}[language=SMT, numbers=none,
caption={SMT type axiom excerpt}]
; Old VC Output
forall ((m T@U) (k Int) (v Int) ) 
  (= (type (q@store m k v)) IntMapType)

; New VC Output
forall ((m T@U) (k Int) (v Int) ) 
  (=> (= (type m) IntMapType)
  (= (type (q@store m k v)) IntMapType))
\end{lstlisting}

The SMT-encoded axiom had been updated to include an additional antecedent (the implication \lstinline[columns=fixed, language=SMT]{=>}).
Using \texttt{git bisect}, we traced this change back to a specific Boogie Pull Request.\footnote{\url{https://github.com/boogie-org/boogie/pull/749}}

Fortunately, fixing the corresponding Isabelle proof was straightforward.
The existing method in Isabelle responsible for resolving goals related to this SMT axiom expected a specific logical structure.
Because the structure had changed via the addition of the antecedent, the rigid proof method failed.
We replaced it with a more flexible proof method, resolving the issue:

\begin{isabellebody}
\begin{isamarkuptext}%
(* In VCExprHelper.thy *)%
\end{isamarkuptext}\isamarkuptrue%
\isakeywordONE{method}\isamarkupfalse%
\ fun{\isacharunderscore}{\kern0pt}output{\isacharunderscore}{\kern0pt}axiom{\isacharprime}{\kern0pt}\ \isakeywordTWO{uses}\ NonEmptyTypes\ {\isacharequal}{\kern0pt}\isanewline
{\isacharparenleft}{\kern0pt}\ {\isacharparenleft}{\kern0pt}simp\ add{\isacharcolon}{\kern0pt}\ Let{\isacharunderscore}{\kern0pt}def{\isacharparenright}{\kern0pt}{\isacharcomma}{\kern0pt}{\isacharparenleft}{\kern0pt}rule\ allI{\isacharparenright}{\kern0pt}{\isacharplus}{\kern0pt}{\isacharcomma}{\kern0pt}\ {\isacharparenleft}{\kern0pt}simp\ split{\isacharcolon}{\kern0pt}\ option{\isachardot}{\kern0pt}split{\isacharparenright}{\kern0pt}{\isacharcomma}{\kern0pt}\isanewline
\ \ {\isacharparenleft}{\kern0pt}rule\ conjI{\isacharparenright}{\kern0pt}{\isacharcomma}{\kern0pt}\ {\isacharparenleft}{\kern0pt}rule\ impI{\isacharparenright}{\kern0pt}{\isacharcomma}{\kern0pt}\isanewline
\ \ {\isacharparenleft}{\kern0pt}rule\ val{\isacharunderscore}{\kern0pt}of{\isacharunderscore}{\kern0pt}closed{\isacharunderscore}{\kern0pt}type{\isacharunderscore}{\kern0pt}correct{\isacharbrackleft}{\kern0pt}OF\ NonEmptyTypes{\isacharbrackright}{\kern0pt}{\isacharparenright}{\kern0pt}{\isacharcomma}{\kern0pt}\isanewline
\ \ assumption{\isacharcomma}{\kern0pt}\ rule\ allI{\isacharcomma}{\kern0pt}\ rule\ impI{\isacharparenright}{\kern0pt}\isanewline
%
\ \isamarkupcmt{Old rigid method}\isanewline
\isanewline
\isakeywordONE{method}\isamarkupfalse%
\ fun{\isacharunderscore}{\kern0pt}output{\isacharunderscore}{\kern0pt}axiom\ \isakeywordTWO{uses}\ NonEmptyTypes\ {\isacharequal}{\kern0pt}\isanewline
{\isacharparenleft}{\kern0pt}\ {\isacharparenleft}{\kern0pt}auto\ simp\ add{\isacharcolon}{\kern0pt}\ Let{\isacharunderscore}{\kern0pt}def\ split{\isacharcolon}{\kern0pt}\ option{\isachardot}{\kern0pt}split\isanewline
\ \ intro{\isacharcolon}{\kern0pt}\ val{\isacharunderscore}{\kern0pt}of{\isacharunderscore}{\kern0pt}closed{\isacharunderscore}{\kern0pt}type{\isacharunderscore}{\kern0pt}correct{\isacharbrackleft}{\kern0pt}OF\ NonEmptyTypes{\isacharbrackright}{\kern0pt}{\isacharparenright}{\kern0pt}\ {\isacharparenright}{\kern0pt}\isanewline
%
\ \isamarkupcmt{New flexible method}
\end{isabellebody}

We verified the usage of this method across the repository\footnote{\lstinline[language=bash, columns=fixed]{grep "fun\_output\_axiom" proofs\_dir -r}}
and confirmed that only the four failing proofs relied upon it.
With this fix in place, all four regressed tests passed successfully,
confirming the successful upgrade of the proof generation module from Boogie v2.11.4 to v3.5.5.
